% =====================================================================
% DRAFT for Chapters/Chapter07.tex (label ch:setup), dissertation repo.
%
% WRITTEN 2026-10-29. Fills the gap identified on that date: the thesis
% states that imbalances are shared in proportion to machine rating
% (ch:setup:multi, ch:setup:multi:dynamic) but never says how the
% scheduled active power of each machine is determined, and never gives
% the voltage-dependent load model as equations. Neither is derivable
% from the text as it stands; a reader would assume the published
% IEEE 39 base-case dispatch, which the builder overwrites in full.
%
% PLACEMENT: as a new \subsection of \cref{ch:setup:benchmarks},
% inserted AFTER the IEEE 39 subsection (ch:setup:multi, i.e. after the
% paragraph anchored at ch:setup:multi:configs) and BEFORE
% \subsection{Dynamic Model of the IEEE 39 Benchmark}
% (ch:setup:multi:dynamic).  That order matters: the dynamic subsection's
% post-outage paragraph and the reference period of
% \cref{ch:setup:scenarios:reference} both refer back to the sharing law
% stated here.
%
% NOTATION: uses the post-2026-10-28 convention (N1--N21): \Wirkleistung,
% \Blindleistung, \Spannung for MW / Mvar / p.u.-on-system-base
% quantities; lower-case \wirkleistung, \blindleistung, \spannung only in
% p.u. of a device rating (capability diagrams of \cref{ch3}).  Machine
% nameplate written as $S_{\mathrm{n},i}$, as in \cref{ch:setup:multi}.
%
% SOURCE OF TRUTH for every statement below (repo qOFO, branch
% minimal/thesis-core, 2026-10-29):
%   network/ieee39/zonal_balancing.py   compute_zonal_gen_dispatch()
%   network/ieee39/build.py:171-183     slack_weight = sn_mva
%   network/ieee39/load_model.py        apply_zip_load_model()
%   network/ieee39/constants.py         LOAD_CONST_FRACTION, PROFILE_MAX,
%                                       DSO_Q_PROFILE_BASE_MVAR, GEN_NAMEPLATE
%   network/ieee39/hv_networks.py:224-292  STS load rows
%   core/actuator_bounds.py:190         compute_generator_q_limits()
%   experiments/runners/multi_tso_dso.py:1833-1846, 3378-3381
%   configs/config.py:1155, 1205, 1216, 1293
% VERIFIED against graphics/setup/reference_period_pf.csv: the law below
% reproduces ne_P = 1240.0 MW at t=0 and trip_gen_P = 562.4 MW at
% t = 29.3 min exactly (see the daily log of 2026-10-29).
% =====================================================================

% ########################################################################
\subsection{Load Model, Time Series, and Active-Power Schedule} \label{ch:setup:operation}
% ########################################################################
The benchmark of \cref{ch:setup:multi} fixes the network; this subsection fixes how it is operated.
Three conventions are needed before any closed-loop result can be reproduced: the voltage dependence of the loads, the time series that drive load and infeed, and the rule that assigns an active power to each synchronous machine.
None of them is a contribution of this work.
\Cref{ch1:sec:scope:dispatch} places active-power dispatch outside the scope, so the schedule below is a declared boundary condition: it must be reproducible and plausible, not economically optimal.

\paragraph{Voltage-dependent load model.}
Every load of the quasi-steady-state plant, at both voltage levels, is an exponential load anchored at a reference voltage $\Spannung_{\mathrm{a}} = \qty{1.03}{\perunit}$,
\begin{equation}
	\Wirkleistung_{l}(t,\Spannung) = \Wirkleistung_{l}^{0}\,\lambda_{\mathrm{p}(l)}(t)\left(\frac{\Spannung}{\Spannung_{\mathrm{a}}}\right)^{\kappa_{\mathrm{p}}},
	\qquad
	\Blindleistung_{l}(t,\Spannung) = \Blindleistung_{l}^{0}\,\lambda_{\mathrm{q}(l)}(t)\left(\frac{\Spannung}{\Spannung_{\mathrm{a}}}\right)^{\kappa_{\mathrm{q}}},
	\label{eq:setup:zip}
\end{equation}
with exponents $\kappa_{\mathrm{p}} = 1$ and $\kappa_{\mathrm{q}} = 2$, nominal powers $\Wirkleistung_{l}^{0}$, $\Blindleistung_{l}^{0}$ and the profile value $\lambda_{\pi}(t)$ of \cref{eq:setup:profile_split} below.
The anchor is a modelling convention, not a physical property: the profile series are voltage-agnostic power scalings, and $\Spannung_{\mathrm{a}}$ only fixes the voltage at which the profile value is served exactly.
It is set to the network voltage setpoint so that the power balance at the initial operating point is the profile balance.
The same exponents are imposed on the \gls{RMS} model of \cref{ch:setup:multi:dynamic}, which is what makes the parity between the two plants a statement about dynamics rather than about load representation.
The loads carry no recovery dynamics, a limitation returned to in \cref{ch:conclusions:limitations}.
Two deliberate exceptions: the reactive load steps of \cref{ch:setup:scenarios:n1} stay constant-power, because they are specified disturbance magnitudes and not physical demand, and the auxiliary buses of \cref{tab:app:dso:buses} exist only to keep a co-located generator from corrupting the load composition of its bus.

\paragraph{Time series.}
Load and infeed follow measured \SI{15}{\minute} profiles of the SimBench dataset \cite{meinecke.2020} over the calendar year \num{2016}, linearly interpolated to the simulation step.
Transmission loads carry the \gls{HV} profile of their area, \texttt{HS4} in areas~1 and~2 and \texttt{HS5} in area~3; sub-transmission loads carry a measured rural profile, wind parks the profiles \texttt{WP7} and \texttt{WP10}, and photovoltaic plants \texttt{PV3}.
Each profile-driven load row is split into a constant and a profile-driven share, as stated in \cref{ch:setup:multi},
\begin{equation}
	\Wirkleistung_{l}^{0}(t) = \Wirkleistung_{l}^{\mathrm{ref}}\left[c + (1-c)\,\frac{\lambda_{\pi}(t)}{\lambda_{\pi}^{\max}}\right],
	\qquad c = \num{0.4},
	\label{eq:setup:profile_split}
\end{equation}
normalised to the annual maximum $\lambda_{\pi}^{\max}$ of the profile, so that the reference load is recovered at the annual peak and the constant share is the asymptotic floor of the affine law.
Sub-transmission active loads follow the same law, distributed over the ten buses with the \numrange{1}{2} weighting of \cref{ch:setup:multi:hv}; their reactive demand, by contrast, has no constant share and follows its profile directly from an aggregate base of \qty{500}{\MVar} per area.
That profile is negative in \SI{85}{\percent} of the year, so the aggregated \gls{HV}/\gls{MV} substations are capacitive most of the time, with a per-area reactive exchange between about \qty{-70}{\MVar} and \qty{+172}{\MVar}.
This is a material boundary condition for every interface result of \cref{pt:CaseStudies} and is the reason the interface reactive flow of \cref{fig:dso:annual:pq} reverses sign over the operating range.

\paragraph{Active-power schedule.}
Reactive power is never scheduled: machines are voltage-controlling buses whose reactive output is the dependent injection of \cref{ch3:sec3:gen}, and \gls{IBG} reactive power is the output of the control layer.
Active power is scheduled open-loop, per control area and per time step, from the profiles alone.
Let $\mathcal{L}_{a}$ and $\mathcal{E}_{a}$ be the loads and the exogenous injectors --- wind and photovoltaic plants --- of area $a$, where both sets include the buses of the \glspl{STS} attached to that area, and let $\mathcal{G}_{a}$ be its dispatchable synchronous machines.
The area residual is
\begin{equation}
	R_{a}(t) = \sum_{l \in \mathcal{L}_{a}} \Wirkleistung_{l}^{0}(t) - \sum_{m \in \mathcal{E}_{a}} \Wirkleistung_{m}^{\mathrm{ref}}\,\lambda_{\pi(m)}(t),
	\label{eq:setup:residual}
\end{equation}
and each machine covers a share of it in proportion to its nameplate rating,
\begin{equation}
	\Wirkleistung_{i}(t) = \min\!\left\{\,R_{a}(t)\,\frac{S_{\mathrm{n},i}}{\sum_{j \in \mathcal{G}_{a}} S_{\mathrm{n},j}},\;\; \alpha\,S_{\mathrm{n},i}\right\},
	\qquad \alpha = \num{0.85},
	\label{eq:setup:dispatch}
\end{equation}
bounded below by zero, so that no technical minimum is imposed.
The cap $\alpha$ keeps every machine inside its capability diagram with usable reactive range: at $\wirkleistung = \alpha$ the stator limit \cref{eq:gen_stator_limit} still admits $\blindleistung \approx \num{0.53}$.
Where the clipping in \cref{eq:setup:dispatch} leaves an area residual uncovered, it is redistributed in proportion to the remaining headroom $\alpha S_{\mathrm{n},i} - \Wirkleistung_{i}$, first inside the area and then, if the area as a whole is capped, across all machines of the system.
The remainder --- the network losses, and any residual beyond $\sum_{i} \alpha S_{\mathrm{n},i}$ --- falls to the distributed slack of \cref{eq:setup:sharing}.

Three consequences carry into the results and are properties of the benchmark, not of the controllers.
The machines are distributed very unevenly over the areas, with \num{12000}, \num{800} and \qty{1500}{\MVA} of dispatchable rating in areas~1 to~3; area~1 holds the \qty{10000}{\MVA} equivalent, which therefore covers about \SI{83}{\percent} of its area residual and, in addition, the losses and the system-wide remainder.
Area~2 faces three exporting \glspl{STS} with a single \qty{800}{\MVA} machine, so its residual is small and its machine lightly loaded, which leaves that machine a wide reactive range.
And because \cref{eq:setup:dispatch} is proportional to rating rather than to marginal cost, the schedule is not a merit-order dispatch and no result of \cref{pt:CaseStudies} may be read as a statement about unit commitment.
What the schedule does guarantee is the property the reactive-power study needs: each area covers its own active residual, so the interface and tie-line flows are set by the area's own load and infeed rather than by an arbitrary allocation.

\paragraph{Imbalance sharing after a disturbance.}
The quasi-steady-state plant has no frequency variable, and the two mechanisms that answer a disturbance must be distinguished.
A change that enters through the profiles --- including a scripted active load ramp --- is answered by re-solving \crefrange{eq:setup:residual}{eq:setup:dispatch} at the next step, that is, by a new schedule.
A contingency is not: generator and line outages and the reactive load steps of \cref{ch:setup:scenarios:n1} never enter the profile frame, so the schedule is untouched and the entire imbalance $\Delta\Wirkleistung$ is taken up by the distributed slack, in proportion to the rating of the machines still in service,
\begin{equation}
	\Delta\Wirkleistung_{i} = \Delta\Wirkleistung \,\frac{S_{\mathrm{n},i}}{\sum_{j \,\in\, \text{in service}} S_{\mathrm{n},j}}.
	\label{eq:setup:sharing}
\end{equation}
\Cref{eq:setup:sharing} is the settled outcome of primary frequency control, and the conditions for that identity are visible in the comparison: a governor with per-unit droop $R_{i}$ shares a disturbance as $\Delta\Wirkleistung_{i} = -(\Delta f / R_{i})\,S_{\mathrm{n},i}$, which reduces to \cref{eq:setup:sharing} if and only if every machine has the same per-unit droop, no governor reaches a limit, the load does not depend on frequency, and no secondary control acts.
The quasi-steady-state plant then reproduces the sharing exactly but at $\Delta f = 0$, and neglects the settled frequency deviation
\begin{equation}
	\frac{\Delta f}{f_{\mathrm{n}}} = -R\,\frac{\Delta\Wirkleistung}{\sum_{j} S_{\mathrm{n},j}}
	\label{eq:setup:df}
\end{equation}
that the dynamic model exhibits.
\Cref{ch:setup:multi:dynamic} checks this representation against the \gls{RMS} model at the operating point after a generator outage, where it matters most, because every post-contingency dispatch of \cref{pt:CaseStudies} starts there.

% =====================================================================
% AMENDMENTS NEEDED ELSEWHERE, once this subsection is inserted
% (2026-10-29; all in the dissertation repo, against origin/main):
%
% 1. Appendices/AppendixB_Benchmarks.tex, IEEE 39 multi-TSO section:
%    "each machine carries a reactive-power envelope of $\pm 0.5\,S_{\mathrm{n}}$"
%    is accurate about the plant data but incomplete and reads as if the
%    envelope were rectangular.  It is the box that the power flow
%    enforces; the CONTROLLER bounds generator reactive power with the
%    P- and V-dependent capability diagram of \cref{ch3:sec3:gen}
%    (stator, rotor and under-excitation limits, re-evaluated at the
%    measured operating point once per iteration,
%    core/actuator_bounds.py:compute_generator_q_limits).  The two agree
%    at the dispatch cap by construction -- at p = 0.85 the stator circle
%    gives q = 0.527 -- and the box is conservative below it.  Suggested
%    clause: "... a reactive-power envelope of $\pm 0.5\,S_{\mathrm{n}}$ in
%    the power flow, which is the stator limit at the dispatch cap
%    $\alpha S_{\mathrm{n}}$ of \cref{eq:setup:dispatch}; the controllers
%    bound the machines with the operating-point-dependent capability
%    diagram of \cref{ch3:sec3:gen} instead."
%
% 2. Chapters/Chapter07.tex, ch:setup:multi, the "Six modifications"
%    paragraph: after "... a weight proportional to its rating", add
%    "(\cref{eq:setup:sharing})" so the clause points at the law instead
%    of standing alone.
%
% 3. Chapters/Chapter07.tex, ch:setup:multi:dynamic: the sentence "A change
%    in active-power balance ... is distributed over the in-service
%    synchronous machines in proportion to their rating" now duplicates
%    \cref{eq:setup:sharing}.  Replace with a pointer and keep only the
%    four conditions, or delete the four conditions here and keep the
%    pointer -- they must not be stated twice.
%
% 4. Chapters/Chapter07.tex, ch:setup:scenarios:reference: "as stated in
%    \cref{ch:setup:multi}" should become "\cref{eq:setup:sharing}".
%
% 5. The settled frequency deviation row of tab:setup:multi:outage is
%    [TBD].  \Cref{eq:setup:df} predicts, for the G 04 outage of the
%    reference period ($\Delta P = \qty{562.4}{\MW}$,
%    $\sum S_{\mathrm{n}} = \qty{13500}{\MVA}$, R = 4 %, f_n = 60 Hz),
%    $\Delta f = \qty{-0.100}{\hertz}$, or -0.098 Hz once the 13.4 MW
%    loss reduction is netted out.  A materially smaller measured value
%    locates load frequency dependence (k_pf, never documented); a larger
%    one locates a governor at its limit.
%
% 6. CHECK before inserting: \cref{eq:gen_stator_limit} and
%    \cref{ch:conclusions:limitations} must exist under those labels in
%    the current build; both were present on origin/main at 5054ae2.
% =====================================================================
